
\documentclass{article}
\usepackage{colortbl}
\usepackage{makecell}
\usepackage{multirow}
\usepackage{supertabular}

\begin{document}

\newcounter{utterance}

\centering \large Interaction Transcript for game `function\_detective', experiment `list\_active\_pair\_checks', episode 74 with gpt{-}5.4.
\vspace{24pt}

{ \footnotesize  \setcounter{utterance}{1}
\setlength{\tabcolsep}{0pt}
\begin{supertabular}{c@{$\;$}|p{.15\linewidth}@{}p{.15\linewidth}p{.15\linewidth}p{.15\linewidth}p{.15\linewidth}p{.15\linewidth}}
    \# & \multicolumn{2}{c}{Player} && \multicolumn{2}{c}{Game Master} \\
    \hline

    \theutterance \stepcounter{utterance}  
    & & & \multicolumn{4}{p{0.6\linewidth}}{
        \cellcolor[rgb]{0.9,0.9,0.9}{
            \makecell[{{p{\linewidth}}}]{
                \texttt{\tiny{[P1$\langle$GM]}}
                \texttt{I have a hidden Python function for you to figure out.} \\
\\ 
\texttt{Your job is to infer what it does and then submit a correct Python implementation.} \\
\\ 
\texttt{Function signature:} \\
\texttt{(items: List{[}int{]}) {-}> int} \\
\\ 
\texttt{You have 15 turns total.} \\
\\ 
\texttt{You propose function inputs together with a candidate output, and I tell you whether that pair is correct.} \\
\\ 
\texttt{You must make your best SOLVE attempt no later than the second{-}to{-}last turn.} \\
\texttt{Do not wait until the very end.} \\
\\ 
\texttt{You can use these actions:} \\
\texttt{{-} JSON format described at the end of the prompt} \\
\texttt{{-} SOLVE: ```python ... ```} \\
\\ 
\texttt{SOLVE: ```python} \\
\texttt{def solution(<value1>):} \\
\texttt{    ...} \\
\texttt{```} \\
\\ 
\texttt{Do not add import statements in your guess.} \\
\texttt{Anything you might need is already available, so just write the function body and return the answer.} \\
\\ 
\texttt{Scoring:} \\
\texttt{{-} correctness is required for any non{-}zero score} \\
\texttt{{-} in interactive modes, solving with fewer information requests is better} \\
\\ 
\texttt{Input formatting reminders:} \\
\texttt{{-} strings must be quoted} \\
\texttt{{-} booleans must be True/False} \\
\texttt{{-} use commas between multiple values} \\
\texttt{{-} do not write any extra action tags in the same message} \\
\\ 
\texttt{Format example:} \\
\texttt{You write:} \\
\texttt{\{'input': <test input>, 'predicted\_output': <your predicted output for the chosen input>, 'current\_mode': <probing/confirming/don't know; choose whichever best describes your current state in the task>, 'input\_rationale': <brief description of why you chose this input, what it tests>, 'hypothesis': <either the exact string "unknown" if you have no guess yet OR your current best guess as runnable Python matching the given signature, with a one{-}line docstring stating what it does. Never write placeholder code, if you cannot commit to a guess, use "unknown". Match the given signature.} \\
\texttt{e.g.} \\
\texttt{def f(x):} \\
\texttt{	"""Returns the absolute value of x."""} \\
\texttt{	return x if x >= 0 else {-}x} \\
\\ 
\texttt{Use "unknown" if you have no guess yet.>, 'confidence\_score': <score between 1{-}100. 100 means highest confidence in your hypothesis>\}} \\
\\ 
\texttt{I answer:} \\
\texttt{OUTPUT: True} \\
\\ 
\texttt{Output only one valid json per turn.} \\
            }
        }
    }
     \\ \\

    \theutterance \stepcounter{utterance}  
    & \multicolumn{4}{p{0.6\linewidth}}{
        \cellcolor[rgb]{0.561,0.737,0.561}{
            \makecell[{{p{\linewidth}}}]{
                \texttt{\tiny{[P1$\rangle$GM]}}
                \texttt{\{"input":{[}2,3,4{]},"predicted\_output":9,"current\_mode":"probing","input\_rationale":"Start by testing the simplest high{-}prior hypothesis: sum of the elements, using an input where common alternatives give different results.","hypothesis":"def solution(items):\textbackslash{}n    \textbackslash{}"\textbackslash{}"\textbackslash{}"Returns the sum of the integers in the list.\textbackslash{}"\textbackslash{}"\textbackslash{}"\textbackslash{}n    return sum(items)","confidence\_score":18\}} \\
            }
        }
    }
    & & \\ \\

    \theutterance \stepcounter{utterance}  
    & & & \multicolumn{4}{p{0.6\linewidth}}{
        \cellcolor[rgb]{0.9,0.9,0.9}{
            \makecell[{{p{\linewidth}}}]{
                \texttt{\tiny{[P1$\langle$GM]}}
                \texttt{OUTPUT: False} \\
            }
        }
    }
     \\ \\

    \theutterance \stepcounter{utterance}  
    & \multicolumn{4}{p{0.6\linewidth}}{
        \cellcolor[rgb]{0.561,0.737,0.561}{
            \makecell[{{p{\linewidth}}}]{
                \texttt{\tiny{[P1$\rangle$GM]}}
                \texttt{\{"input":{[}5,{-}1,2{]},"predicted\_output":5,"current\_mode":"probing","input\_rationale":"Test the common hypothesis that the function returns the maximum element, using mixed{-}sign values to separate max from sum, min, len, and product.","hypothesis":"def solution(items):\textbackslash{}n    \textbackslash{}"\textbackslash{}"\textbackslash{}"Returns the maximum integer in the list.\textbackslash{}"\textbackslash{}"\textbackslash{}"\textbackslash{}n    return max(items)","confidence\_score":16\}} \\
            }
        }
    }
    & & \\ \\

    \theutterance \stepcounter{utterance}  
    & & & \multicolumn{4}{p{0.6\linewidth}}{
        \cellcolor[rgb]{0.9,0.9,0.9}{
            \makecell[{{p{\linewidth}}}]{
                \texttt{\tiny{[P1$\langle$GM]}}
                \texttt{OUTPUT: False} \\
            }
        }
    }
     \\ \\

    \theutterance \stepcounter{utterance}  
    & \multicolumn{4}{p{0.6\linewidth}}{
        \cellcolor[rgb]{0.561,0.737,0.561}{
            \makecell[{{p{\linewidth}}}]{
                \texttt{\tiny{[P1$\rangle$GM]}}
                \texttt{\{"input":{[}7,8,9,10{]},"predicted\_output":4,"current\_mode":"probing","input\_rationale":"Check another very common list{-}to{-}int behavior: returning the length. This input avoids accidental matches with max, min, sum, or range.","hypothesis":"def solution(items):\textbackslash{}n    \textbackslash{}"\textbackslash{}"\textbackslash{}"Returns the number of elements in the list.\textbackslash{}"\textbackslash{}"\textbackslash{}"\textbackslash{}n    return len(items)","confidence\_score":15\}} \\
            }
        }
    }
    & & \\ \\

    \theutterance \stepcounter{utterance}  
    & & & \multicolumn{4}{p{0.6\linewidth}}{
        \cellcolor[rgb]{0.9,0.9,0.9}{
            \makecell[{{p{\linewidth}}}]{
                \texttt{\tiny{[P1$\langle$GM]}}
                \texttt{OUTPUT: False} \\
            }
        }
    }
     \\ \\

    \theutterance \stepcounter{utterance}  
    & \multicolumn{4}{p{0.6\linewidth}}{
        \cellcolor[rgb]{0.561,0.737,0.561}{
            \makecell[{{p{\linewidth}}}]{
                \texttt{\tiny{[P1$\rangle$GM]}}
                \texttt{\{"input":{[}2,2,2{]},"predicted\_output":2,"current\_mode":"probing","input\_rationale":"Classify whether the function tends to return a representative element/value from the list (min/max/median/gcd/etc.) rather than a count or arithmetic aggregate like sum or length.","hypothesis":"unknown","confidence\_score":10\}} \\
            }
        }
    }
    & & \\ \\

    \theutterance \stepcounter{utterance}  
    & & & \multicolumn{4}{p{0.6\linewidth}}{
        \cellcolor[rgb]{0.9,0.9,0.9}{
            \makecell[{{p{\linewidth}}}]{
                \texttt{\tiny{[P1$\langle$GM]}}
                \texttt{OUTPUT: False} \\
            }
        }
    }
     \\ \\

    \theutterance \stepcounter{utterance}  
    & \multicolumn{4}{p{0.6\linewidth}}{
        \cellcolor[rgb]{0.561,0.737,0.561}{
            \makecell[{{p{\linewidth}}}]{
                \texttt{\tiny{[P1$\rangle$GM]}}
                \texttt{\{"input":{[}{-}3,5{]},"predicted\_output":8,"current\_mode":"probing","input\_rationale":"Test the hypothesis that the function returns max(items) {-} min(items). This input makes that value distinctive from sum, len, max, min, product, and counts.","hypothesis":"def solution(items):\textbackslash{}n    \textbackslash{}"\textbackslash{}"\textbackslash{}"Returns the difference between the largest and smallest integers in the list.\textbackslash{}"\textbackslash{}"\textbackslash{}"\textbackslash{}n    return max(items) {-} min(items)","confidence\_score":22\}} \\
            }
        }
    }
    & & \\ \\

    \theutterance \stepcounter{utterance}  
    & & & \multicolumn{4}{p{0.6\linewidth}}{
        \cellcolor[rgb]{0.9,0.9,0.9}{
            \makecell[{{p{\linewidth}}}]{
                \texttt{\tiny{[P1$\langle$GM]}}
                \texttt{OUTPUT: False} \\
            }
        }
    }
     \\ \\

    \theutterance \stepcounter{utterance}  
    & \multicolumn{4}{p{0.6\linewidth}}{
        \cellcolor[rgb]{0.561,0.737,0.561}{
            \makecell[{{p{\linewidth}}}]{
                \texttt{\tiny{[P1$\rangle$GM]}}
                \texttt{\{"input":{[}0,0,0{]},"predicted\_output":0,"current\_mode":"probing","input\_rationale":"Many list{-}to{-}int functions collapse an all{-}zero input to 0 (sum/product/range/xor/max{-}subarray/count of positives, etc.). This is a high{-}probability attempt to get one confirmed datapoint.","hypothesis":"unknown","confidence\_score":20\}} \\
            }
        }
    }
    & & \\ \\

    \theutterance \stepcounter{utterance}  
    & & & \multicolumn{4}{p{0.6\linewidth}}{
        \cellcolor[rgb]{0.9,0.9,0.9}{
            \makecell[{{p{\linewidth}}}]{
                \texttt{\tiny{[P1$\langle$GM]}}
                \texttt{OUTPUT: False} \\
            }
        }
    }
     \\ \\

    \theutterance \stepcounter{utterance}  
    & \multicolumn{4}{p{0.6\linewidth}}{
        \cellcolor[rgb]{0.561,0.737,0.561}{
            \makecell[{{p{\linewidth}}}]{
                \texttt{\tiny{[P1$\rangle$GM]}}
                \texttt{\{"input":{[}1,1,1{]},"predicted\_output":3,"current\_mode":"probing","input\_rationale":"Use an input where many common functions coincide at 3 (sum, length, count of positives, product, etc.) to maximize the chance of obtaining a first confirmed datapoint.","hypothesis":"unknown","confidence\_score":24\}} \\
            }
        }
    }
    & & \\ \\

    \theutterance \stepcounter{utterance}  
    & & & \multicolumn{4}{p{0.6\linewidth}}{
        \cellcolor[rgb]{0.9,0.9,0.9}{
            \makecell[{{p{\linewidth}}}]{
                \texttt{\tiny{[P1$\langle$GM]}}
                \texttt{OUTPUT: False} \\
            }
        }
    }
     \\ \\

    \theutterance \stepcounter{utterance}  
    & \multicolumn{4}{p{0.6\linewidth}}{
        \cellcolor[rgb]{0.561,0.737,0.561}{
            \makecell[{{p{\linewidth}}}]{
                \texttt{\tiny{[P1$\rangle$GM]}}
                \texttt{\{"input":{[}0,0,0{]},"predicted\_output":1,"current\_mode":"probing","input\_rationale":"All{-}equal zeros are a strong separator: many arithmetic/counting functions return 0, while many predicate/index{-}style functions return 1. I want a first positive anchor.","hypothesis":"unknown","confidence\_score":29\}} \\
            }
        }
    }
    & & \\ \\

    \theutterance \stepcounter{utterance}  
    & & & \multicolumn{4}{p{0.6\linewidth}}{
        \cellcolor[rgb]{0.9,0.9,0.9}{
            \makecell[{{p{\linewidth}}}]{
                \texttt{\tiny{[P1$\langle$GM]}}
                \texttt{OUTPUT: False} \\
            }
        }
    }
     \\ \\

    \theutterance \stepcounter{utterance}  
    & \multicolumn{4}{p{0.6\linewidth}}{
        \cellcolor[rgb]{0.561,0.737,0.561}{
            \makecell[{{p{\linewidth}}}]{
                \texttt{\tiny{[P1$\rangle$GM]}}
                \texttt{\{"input":{[}0,0,0{]},"predicted\_output":3,"current\_mode":"probing","input\_rationale":"We know the output for this input is neither 0 nor 1. Testing 3 is a strong way to distinguish common count{-}based behaviors like number of even elements or number of zeros.","hypothesis":"unknown","confidence\_score":35\}} \\
            }
        }
    }
    & & \\ \\

    \theutterance \stepcounter{utterance}  
    & & & \multicolumn{4}{p{0.6\linewidth}}{
        \cellcolor[rgb]{0.9,0.9,0.9}{
            \makecell[{{p{\linewidth}}}]{
                \texttt{\tiny{[P1$\langle$GM]}}
                \texttt{OUTPUT: False} \\
            }
        }
    }
     \\ \\

    \theutterance \stepcounter{utterance}  
    & \multicolumn{4}{p{0.6\linewidth}}{
        \cellcolor[rgb]{0.561,0.737,0.561}{
            \makecell[{{p{\linewidth}}}]{
                \texttt{\tiny{[P1$\rangle$GM]}}
                \texttt{\{"input":{[}1,2{]},"predicted\_output":2,"current\_mode":"probing","input\_rationale":"For a 2{-}element increasing list, output 2 is shared by many common hypotheses (length, max, product, lcm, number of distinct elements). A True here would anchor the function quickly.","hypothesis":"unknown","confidence\_score":33\}} \\
            }
        }
    }
    & & \\ \\

    \theutterance \stepcounter{utterance}  
    & & & \multicolumn{4}{p{0.6\linewidth}}{
        \cellcolor[rgb]{0.9,0.9,0.9}{
            \makecell[{{p{\linewidth}}}]{
                \texttt{\tiny{[P1$\langle$GM]}}
                \texttt{OUTPUT: False} \\
            }
        }
    }
     \\ \\

    \theutterance \stepcounter{utterance}  
    & \multicolumn{4}{p{0.6\linewidth}}{
        \cellcolor[rgb]{0.561,0.737,0.561}{
            \makecell[{{p{\linewidth}}}]{
                \texttt{\tiny{[P1$\rangle$GM]}}
                \texttt{\{"input":{[}1,{-}1,2,{-}2{]},"predicted\_output":3,"current\_mode":"probing","input\_rationale":"Test the hypothesis that the function counts zero{-}sum contiguous subarrays. This input has exactly 3 such subarrays, while many other plausible subarray{-}count functions give different values.","hypothesis":"def solution(items):\textbackslash{}n    \textbackslash{}"\textbackslash{}"\textbackslash{}"Returns the number of contiguous subarrays whose sum is zero.\textbackslash{}"\textbackslash{}"\textbackslash{}"\textbackslash{}n    total = 0\textbackslash{}n    pref = 0\textbackslash{}n    counts = \{0: 1\}\textbackslash{}n    for x in items:\textbackslash{}n        pref += x\textbackslash{}n        total += counts.get(pref, 0)\textbackslash{}n        counts{[}pref{]} = counts.get(pref, 0) + 1\textbackslash{}n    return total","confidence\_score":41\}} \\
            }
        }
    }
    & & \\ \\

    \theutterance \stepcounter{utterance}  
    & & & \multicolumn{4}{p{0.6\linewidth}}{
        \cellcolor[rgb]{0.9,0.9,0.9}{
            \makecell[{{p{\linewidth}}}]{
                \texttt{\tiny{[P1$\langle$GM]}}
                \texttt{OUTPUT: False} \\
            }
        }
    }
     \\ \\

    \theutterance \stepcounter{utterance}  
    & \multicolumn{4}{p{0.6\linewidth}}{
        \cellcolor[rgb]{0.561,0.737,0.561}{
            \makecell[{{p{\linewidth}}}]{
                \texttt{\tiny{[P1$\rangle$GM]}}
                \texttt{\{"input":{[}0,0,0{]},"predicted\_output":6,"current\_mode":"probing","input\_rationale":"A lot of plausible list{-}to{-}int functions that survived earlier eliminations map three zeros to 6: total number of contiguous subarrays, count of zero{-}sum subarrays, count of even{-}sum subarrays. This is a high{-}information check.","hypothesis":"unknown","confidence\_score":44\}} \\
            }
        }
    }
    & & \\ \\

    \theutterance \stepcounter{utterance}  
    & & & \multicolumn{4}{p{0.6\linewidth}}{
        \cellcolor[rgb]{0.9,0.9,0.9}{
            \makecell[{{p{\linewidth}}}]{
                \texttt{\tiny{[P1$\langle$GM]}}
                \texttt{OUTPUT: False} \\
            }
        }
    }
     \\ \\

    \theutterance \stepcounter{utterance}  
    & \multicolumn{4}{p{0.6\linewidth}}{
        \cellcolor[rgb]{0.561,0.737,0.561}{
            \makecell[{{p{\linewidth}}}]{
                \texttt{\tiny{[P1$\rangle$GM]}}
                \texttt{\{"input":{[}1,{-}1,2,{-}2{]},"predicted\_output":4,"current\_mode":"confirming","input\_rationale":"Test the focused hypothesis that the function counts subsets whose sum is zero, including the empty subset. For this input, those subsets are \{\}, \{1,{-}1\}, \{2,{-}2\}, and \{1,{-}1,2,{-}2\}, giving 4.","hypothesis":"def solution(items):\textbackslash{}n    \textbackslash{}"\textbackslash{}"\textbackslash{}"Returns the number of subsets whose sum is zero, including the empty subset.\textbackslash{}"\textbackslash{}"\textbackslash{}"\textbackslash{}n    n = len(items)\textbackslash{}n    total = 0\textbackslash{}n    for mask in range(1 << n):\textbackslash{}n        s = 0\textbackslash{}n        for i in range(n):\textbackslash{}n            if mask \& (1 << i):\textbackslash{}n                s += items{[}i{]}\textbackslash{}n        if s == 0:\textbackslash{}n            total += 1\textbackslash{}n    return total","confidence\_score":57\}} \\
            }
        }
    }
    & & \\ \\

    \theutterance \stepcounter{utterance}  
    & & & \multicolumn{4}{p{0.6\linewidth}}{
        \cellcolor[rgb]{0.9,0.9,0.9}{
            \makecell[{{p{\linewidth}}}]{
                \texttt{\tiny{[P1$\langle$GM]}}
                \texttt{OUTPUT: False} \\
            }
        }
    }
     \\ \\

    \theutterance \stepcounter{utterance}  
    & \multicolumn{4}{p{0.6\linewidth}}{
        \cellcolor[rgb]{0.561,0.737,0.561}{
            \makecell[{{p{\linewidth}}}]{
                \texttt{\tiny{[P1$\rangle$GM]}}
                \texttt{\{"input":{[}2,1{]},"predicted\_output":2,"current\_mode":"confirming","input\_rationale":"This distinguishes a strong surviving hypothesis: counting non{-}empty non{-}decreasing subsequences. For {[}2,1{]}, only the singletons qualify, so the count is 2. It also separates that idea from length{-}only subset counts (which would give 3 or 4).","hypothesis":"def solution(items):\textbackslash{}n    \textbackslash{}"\textbackslash{}"\textbackslash{}"Returns the number of non{-}empty non{-}decreasing subsequences in the list.\textbackslash{}"\textbackslash{}"\textbackslash{}"\textbackslash{}n    n = len(items)\textbackslash{}n    total = 0\textbackslash{}n    for mask in range(1, 1 << n):\textbackslash{}n        last = None\textbackslash{}n        ok = True\textbackslash{}n        for i in range(n):\textbackslash{}n            if mask \& (1 << i):\textbackslash{}n                if last is not None and items{[}i{]} < last:\textbackslash{}n                    ok = False\textbackslash{}n                    break\textbackslash{}n                last = items{[}i{]}\textbackslash{}n        if ok:\textbackslash{}n            total += 1\textbackslash{}n    return total","confidence\_score":52\}} \\
            }
        }
    }
    & & \\ \\

    \theutterance \stepcounter{utterance}  
    & & & \multicolumn{4}{p{0.6\linewidth}}{
        \cellcolor[rgb]{0.9,0.9,0.9}{
            \makecell[{{p{\linewidth}}}]{
                \texttt{\tiny{[P1$\langle$GM]}}
                \texttt{OUTPUT: False} \\
            }
        }
    }
     \\ \\

    \theutterance \stepcounter{utterance}  
    & \multicolumn{4}{p{0.6\linewidth}}{
        \cellcolor[rgb]{0.561,0.737,0.561}{
            \makecell[{{p{\linewidth}}}]{
                \texttt{\tiny{[P1$\rangle$GM]}}
                \texttt{SOLVE: ```python} \\
\texttt{def solution(items):} \\
\texttt{    """Returns the number of non{-}empty subsets/subsequences of the list."""} \\
\texttt{    return (1 << len(items)) {-} 1} \\
\texttt{```} \\
            }
        }
    }
    & & \\ \\

    \theutterance \stepcounter{utterance}  
    & & & \multicolumn{2}{p{0.3\linewidth}}{
        \cellcolor[rgb]{0.9,0.9,0.9}{
            \makecell[{{p{\linewidth}}}]{
                \texttt{\tiny{[GM$|$GM]}}
                \texttt{Game Verdict: LOSS} \\
            }
        }
    }
    & & \\ \\

    \theutterance \stepcounter{utterance}  
    & & & \multicolumn{2}{p{0.3\linewidth}}{
        \cellcolor[rgb]{0.9,0.9,0.9}{
            \makecell[{{p{\linewidth}}}]{
                \texttt{\tiny{[GM$|$GM]}}
                \texttt{Incorrect.} \\
\texttt{Failed Cases:} \\
\texttt{Input: {[}{[}5, {-}5, 5{]}{]} | Expected: 2 | Got: 7} \\
\texttt{Input: {[}{[}1, 2, 3{]}{]} | Expected: 0 | Got: 7} \\
\texttt{Input: {[}{[}0{]}{]} | Expected: 0 | Got: 1} \\
            }
        }
    }
    & & \\ \\

    \theutterance \stepcounter{utterance}  
    & & & \multicolumn{2}{p{0.3\linewidth}}{
        \cellcolor[rgb]{0.9,0.9,0.9}{
            \makecell[{{p{\linewidth}}}]{
                \texttt{\tiny{[GM$|$GM]}}
                \texttt{def count\_adjacent\_opposite\_pairs(items: List{[}int{]}) {-}> int:} \\
\texttt{    return sum(1 for i in range(len(items) {-} 1) if items{[}i{]} == {-}items{[}i + 1{]})} \\
            }
        }
    }
    & & \\ \\

    \theutterance \stepcounter{utterance}  
    & & & \multicolumn{2}{p{0.3\linewidth}}{
        \cellcolor[rgb]{0.9,0.9,0.9}{
            \makecell[{{p{\linewidth}}}]{
                \texttt{\tiny{[GM$|$GM]}}
                \texttt{def solution(items):} \\
\texttt{    """Returns the number of non{-}empty subsets/subsequences of the list."""} \\
\texttt{    return (1 << len(items)) {-} 1} \\
            }
        }
    }
    & & \\ \\

    \theutterance \stepcounter{utterance}  
    & & & \multicolumn{2}{p{0.3\linewidth}}{
        \cellcolor[rgb]{0.9,0.9,0.9}{
            \makecell[{{p{\linewidth}}}]{
                \texttt{\tiny{[GM$|$GM]}}
                \texttt{Mode: active\_pair\_checks} \\
\texttt{Accuracy: 0.010} \\
\texttt{Information requests: 14/15} \\
\texttt{Revealed passive examples: 0} \\
\texttt{Efficiency: 0.133} \\
\texttt{Slack: 1} \\
\texttt{Observed{-}pair consistency: 0.000} \\
\texttt{Membership balance: 1 true / 14 false} \\
\texttt{Parse errors: 0} \\
\texttt{Runtime errors: 0} \\
            }
        }
    }
    & & \\ \\

\end{supertabular}
}

\end{document}
